% BAN TIENG VIET cua 02_threat_attack.tex (2026-10-08). So lieu giu nguyen.
\section{Mô hình đe doạ}
\label{sec:threat}

Một nhật ký tưới khai báo là một ánh xạ hữu hạn
\begin{equation}
L=\{(t_i,u_i)\}_{i=1}^{|L|},\qquad u_i>0,\ t_i\in[0,T),
\label{eq:log}
\end{equation}
trong đó $u_i$ là độ sâu khai báo tính bằng milimét và $t_i$ là chỉ số giờ.
Ba đại lượng độc lập có thể bị làm giả, và mỗi tầng phòng thủ quan sát một tập
con khác nhau của chúng:
\begin{equation}
\underbrace{\textstyle\sum_i u_i}_{\text{tổng lượng}}\,,\qquad
\underbrace{\{u_i\}}_{\text{từng độ sâu}}\,,\qquad
\underbrace{\{t_i\}}_{\text{dấu thời gian}}\,.
\label{eq:three}
\end{equation}
Bảng~S1 của Phụ lục liệt kê các lớp gian lận tương ứng. Đối thủ là
\emph{duy lý}: nó tối đa hoá số pha khô được cấp tín chỉ với điều kiện không
bị đánh dấu, và nó biết bộ phát hiện, vì bộ phát hiện đã được công bố; đó là
nghĩa mà theo đó $\varepsilon$ trong \eqref{eq:v4} bên dưới không phải là bí
mật.

Gọi $M$ là lượng nước thật sự đã bơm, đo tại trạm của hợp tác xã, và
$U(L)=\sum_i u_i$ là tổng lượng khai báo; thống kê của tầng~1 là
\begin{equation}
G=M-U(L),
\label{eq:gap}
\end{equation}
và đại lượng liên quan đến tín chỉ thì đếm \emph{số pha khô}. Viết
$\mathcal{A}(\cdot)$ cho tập các khoảng tối đại $[a,b)$ của một quỹ đạo phát
lại mà trên đó độ cạn kiệt luôn ở mức từ $d_{\mathrm{dry}}$ trở lên và không
có khai báo tưới nào, kéo dài ít nhất $72$~giờ --- mức tối thiểu ba ngày của
Quyết định 4801~\cite{qd4801}:
\begin{equation}
\begin{aligned}
\mathcal{A}\bigl(\{\hat d_t\}\bigr)=\Bigl\{[a,b):\ &\min_{t\in[a,b)}
\hat d_t\ge d_{\mathrm{dry}},\\
&u_t=0\ \forall t\in[a,b),\\
&b-a\ge 72\,\mathrm{h}\Bigr\}.
\end{aligned}
\label{eq:awd}
\end{equation}
Số pha được cấp tín chỉ khi đó là
\begin{equation}
n_{\mathrm{dry}}(L)=\bigl|\mathcal{A}\bigl(\mathcal{R}(L)\bigr)\bigr|,
\label{eq:ndry}
\end{equation}
với $d_{\mathrm{dry}}=0{,}45$ trong toàn bộ bài. Doanh thu tỉ lệ với
$n_{\mathrm{dry}}$ chứ không tỉ lệ với $U(L)$ --- chính sự phân biệt này làm
cho tấn công ở Mục~\ref{sec:timeshift} sinh lợi, và nó vô hình trong một
cách phát biểu coi tổng lượng nước là mục tiêu.

\section{Các bộ lọc dựa trên vật lý và đúng điểm mù của chúng}
\label{sec:filters}

Phát lại $L$ dưới dữ liệu ERA5 công khai cho ra $(\hat w_t,\hat d_t)$ (công
trình đồng hành~\cite{twingate2026} trình bày đầy đủ cân bằng này); có ba bộ
lọc được áp dụng. \emph{$V_1$, bão hoà}: một khai báo phát ra khi tầng rễ đã
ở sức chứa ẩm đồng ruộng thì không thể được lưu trữ:
\begin{equation}
V_1(L)=\{\,i:\ u_i>0\ \wedge\ \hat w_{t_i}^{-}\ \ge\ FC-\varepsilon\,\}.
\label{eq:v1}
\end{equation}

\emph{$V_4$, vượt sức chứa}: một khai báo lớn hơn phần thiếu hụt sẵn có tại
giờ đó hẳn đã chảy tràn hoặc thấm mất, nên nó không ``được áp dụng'' theo
nghĩa mà quy chuẩn đòi hỏi:
\begin{equation}
V_4(L)=\{\,i:\ u_i\ >\ \underbrace{(FC-\hat w_{t_i}^{-})}_{\text{phần thiếu hụt}}
\ +\ \varepsilon\,\}.
\label{eq:v4}
\end{equation}

\emph{$V_5$, giới hạn thấm.} Không phụ thuộc trạng thái, một lần tưới kéo dài
$\tau_i$ không thể cấp nhiều hơn lượng đất tiếp nhận được:
\begin{equation}
u_i\ \le\ k_{\mathrm{inf}}\,\tau_i+p_{\max},
\label{eq:v5}
\end{equation}
với $k_{\mathrm{inf}}=10$~mm/h và $p_{\max}=100$~mm, lấy ở đầu rộng của dải
giá trị đã công bố cho ruộng lúa làm dầm (giới hạn $110$~mm tại
$\tau_i=1$~giờ), để một nông dân ghi hai lần tưới liền kề trong cùng một giờ
không bị kết tội.

Bộ lọc thứ tư, về độ cạn kiệt xuống dưới điểm héo, đã có trong phiên bản công
bố của lớp này. Nó không thể kích hoạt: xem Mục~\ref{sec:v2reach}.

\begin{proposition}[Điểm mù chính xác của các phòng thủ tổng hợp]
\label{prop:blind}
Giả sử $L$ và $L'$ có cùng đa tập hợp các độ sâu và cùng tổng lượng, tức
$\{u_i\}=\{u'_j\}$ theo nghĩa đa tập hợp. Khi đó
\begin{equation}
G(L)=G(L')\qquad\text{và}\qquad |V_5(L)|=|V_5(L')|.
\end{equation}
\end{proposition}

\begin{proof}
Xem Phụ lục~S11; cả hai đại lượng đều là hàm đối xứng của đa tập hợp độ sâu
và không bao giờ phụ thuộc $t_i$.
\end{proof}

Mệnh đề~\ref{prop:blind} gói trọn toàn bộ khó khăn trong một dòng: hai phòng
thủ không phụ thuộc việc mô hình có đúng hay không --- đồng hồ đo bơm và giới
hạn thấm --- đều bất biến \emph{chính xác} dưới phép biến đổi mà đối thủ cần.
Còn lại $V_1$ và $V_4$, vốn phụ thuộc thứ tự qua $\hat w_{t_i}^{-}$; hoá ra
chúng không giúp được gì.

\section{Tấn công dịch thời điểm}
\label{sec:timeshift}

\begin{proposition}[Lợi ích tín chỉ và sự lọt lưới là cùng một hành động]
\label{prop:coincide}
Giả sử một khai báo tại giờ $t$ có phần thiếu hụt sẵn có
$c_t=FC-\hat w_t^{-}$. Giữa hai khai báo liên tiếp, thùng chứa mất đi
$\Delta_t=\sum_{s}(\mathit{ETc}_s+\mathit{DP}_s-P_s^{+})\ge 0$, nên việc hoãn
một khai báo thêm $\delta$ giờ làm tăng phần thiếu hụt của nó đúng bằng lượng
nước đã thoát đi trong các giờ đó:
\begin{equation}
c_{t+\delta}\ \ge\ c_t+\min\!\bigl(\delta\cdot\overline{\Delta},\ FC-\hat w_t^{-}\bigr).
\label{eq:delay}
\end{equation}
Do đó phép biến đổi dời các khai báo \emph{ra khỏi} một khoảng khô và đặt lên
biên của khoảng đó đồng thời (i)~làm tăng $n_{\mathrm{dry}}$, vì khoảng đó
không còn bị ngắt quãng, và (ii)~làm tăng $c_t$ tại mỗi khai báo bị dời, điều
mà theo \eqref{eq:v4} khiến $V_4$ \emph{ít} có khả năng kích hoạt hơn. Đối
thủ không phải đánh đổi điều gì: gradient của lợi nhuận và gradient của sự
không thể phát hiện cùng chỉ một hướng.
\end{proposition}

\begin{proof}
Xem Phụ lục~S11. Nói ngắn gọn: (i) đúng theo định nghĩa, vì $\mathcal{A}$ đòi
hỏi không có khai báo nào bên trong khoảng; (ii) suy ra được vì hoãn một khai
báo để thùng chứa tiếp tục thoát nước thêm $\delta$ giờ mà không có dòng vào,
nên phần thiếu hụt sẵn có chỉ có thể tăng trong khi $u_i$ không đổi.
\end{proof}

Đó là lý do tấn công này không thể bị bịt bằng thuỷ văn tốt hơn: bất kỳ bộ
phát hiện nào có phán quyết là hàm của quỹ đạo phát lại đều phải tích phân
đúng lượng nước thoát đã làm cho khai báo bị hoãn trông có vẻ hợp pháp. Bản
thực nghiệm của tấn công là một thuật toán leo đồi trên các độ lệch giờ
nguyên trong $[-36,36]$, chỉ chấp nhận một bước dời nếu tổng lượng được bảo
toàn từng bit, mọi độ sâu vẫn ở dưới \eqref{eq:v5}, và $n_{\mathrm{dry}}$
tăng nghiêm ngặt. Tấn công không ăn cắp nước --- thứ bị ăn cắp là
\emph{lịch sử} (Phụ lục~S19).
